\documentclass{article}

\begin{document}

En este documento haremos un pseudocódigo para calcular el tiempo en que tarda un cuerpo
en alcanzar una altura h, tendremos en cuenta la altura inicial $h_0$ y la velocidad inicial $v_0$.

Recordemos que la fórmula de la altura es:
\[h = h_0 + v_0t - \frac{1}{2}gt^2\]

Que al reordenarla de la forma $ax^2 + bx + c$ nos queda:
\[\frac{1}{2}gt^2 - v_0t + (h-h_0) = 0\]

De la cual podemos obtener t mediante la fórmula cuadrática:
\[t = \frac{v_0 \pm \sqrt{v_0^2 - 2g(h-h_0)}}{g}\]

\section{pseudocódigo}

1. Iniciar

2. Definir h preguntándole al usuario (un negativo es perfectamente válido)

3. Definir $h_0$ preguntándole al usuario (un negativo es perfectamente válido)

4. Definir $v_0$ preguntándole al usuario, sesiorándonos de que sea positiva (pues si no, tendríamos que calcular rebotes)

5. Definir g = -9.81

6. Calcular y definir $v_0^2$

7. Calcular y definir $2g(h-h_0)$

8. Haremos un si y si no

8.1. Si $v_0^2 < 2g(h-h_0)$ indicarle al usuario que el objeto jamás alcanzará la altura deseada

8.2. Si no, calcular el tiempo que le tarda al objeto alcanzar la altura deseada

9. Con la altura calculada, mostrar ambos tiempos, indicándole al usuario que debe sesiorarse de cuáles son válidos.

\end{document}